\documentclass[11pt]{article}

% --- Packages ---
\usepackage[margin=1in]{geometry}
\usepackage{setspace}
\onehalfspacing
\setlength{\parskip}{0.15em}
\usepackage{amsmath,amssymb}
\usepackage{booktabs}
\usepackage{graphicx}
\usepackage[hidelinks]{hyperref}
\usepackage[numbers]{natbib}
\usepackage{caption}
\usepackage{enumitem}
\usepackage{float}
\usepackage{url}
\usepackage{titlesec}
\titlespacing*{\section}{0pt}{0.45em}{0.15em}
\titlespacing*{\subsection}{0pt}{0.25em}{0.08em}
\titlespacing*{\paragraph}{0pt}{0.25em}{0.6em}


% --- Title ---
\title{Capacity Portfolio Evaluation Under Multi-Dimensional Uncertainty:\\A Stochastic Simulation Study}
\author{Omar Mourssi --- 1008311961}
\date{April 15, 2026}

\begin{document}

\maketitle

%==========================================================================
% SECTION 1: INTRODUCTION
%==========================================================================
\section{Introduction}\label{sec:intro}

Long-term electricity capacity planning involves selecting a resource mix that balances cost and reliability. The IESO projects a 75\% increase in Ontario's electricity demand by 2050 \citep{IESO2025}; however, as the generation fleet transitions toward variable renewables and battery storage, planning exercises that evaluate only a limited set of fixed scenarios may fail to identify the point at which uncertainty alters the preferred resource mix. The extent to which this occurs depends on the stringency of the capacity budget. This report seeks to quantify that dependence within a stylized modelling framework.

A one-zone system is simulated chronologically over 8{,}760 hours with inputs calibrated from Ontario empirical data, and the total installed capacity is fixed at 12~GW. This 12~GW figure does not represent an Ontario adequacy target; rather, it serves as a budget-constrained stress scenario in which total capacity is insufficient for full adequacy even in expectation. Consequently, the trade-off between firm generation and variable resources is exposed to tail risk. Within this budget and the original 16-portfolio simplex (firm share capped at 70\%), stochastic evaluation selects P04 (70/20/10) as the in-budget winner, while mean-input evaluation selects P03 (60/30/10). The Deterministic Benchmark Gap between them is \$3{,}968M per year on the top-4 set. When the same four portfolios are re-evaluated at 15.8~GW and 18~GW, both methods converge on P03, and the DBG closes within the indifference zone. A separate boundary probe at 12~GW shows that relaxing the 70\% firm cap reveals a still cheaper mix: P$_{b3}$ (85/10/05) beats P04 by \$746M. P04 is therefore the in-budget winner within the original constrained design, not an absolute recommendation, and the value of uncertainty-aware evaluation in this model is concentrated where the capacity budget binds.

Recent literature addresses uncertainty in capacity planning directly.\citet{Roald2023} review optimization under uncertainty in power systems, and \citet{Stover2023} argue that deterministic reserve-margin rules are insufficient for high-renewable systems.\citet{Leibowicz2024} demonstrate that operational detail in Monte Carlo adequacy simulation influences which resources are identified as adequate, while \citet{Wang2024} find that temporal structure is critical for the adequacy contribution of renewables and storage. This motivates the use of NORTA AR(1) for wind input in the present study, rather than independent sampling.\citet{ParkBaldick2016} serves as a motivating reference: they solve a multi-year stochastic capacity-expansion program and find that uncertainty affects optimal build timing and resource mix. The present study addresses a related question in a different context, specifically a fixed-budget evaluation of a discretized candidate set using Kim--Nelson (KN) ranking and selection \citep{KimNelson2001} with Common Random Numbers (CRN) and Bonferroni-corrected confidence intervals, rather than a multi-stage stochastic program. While the two approaches agree in direction within this stylized setting, they are not directly comparable in methodology or decision problem.

Portfolios are generated by discretizing the feasible simplex at 10\% increments subject to $f^{\text{firm}} \in [0.20, 0.70]$, $f^{\text{renew}} \in [0.10, 0.60]$, $f^{\text{stor}} \in [0.10, 0.30]$, giving $k = 16$ candidates. Inputs are calibrated from IESO records and the CODERS dataset \citep{CODERS2025}.

The remainder of this report is structured as follows. Section~\ref{sec:problem} defines the problem and performance measure. Section~\ref{sec:model} describes the simulation model. Section~\ref{sec:design} presents the experimental design. Section~\ref{sec:results} reports the results. Section~\ref{sec:conclusion} discusses limitations and directions for future work.

%==========================================================================
% SECTION 2: PROBLEM DESCRIPTION
%==========================================================================
\section{Problem Description}\label{sec:problem}

\subsection{Portfolio Generation}

The system comprises three technology classes: firm generation (nuclear, flexible firm, mid-merit gas, gas peakers), variable renewables (wind, solar), and battery storage. Total installed capacity is fixed at $C_{\text{total}} = 12{,}000$~MW, derived from 8,000~MW peak demand and a 50\% reserve margin. Each portfolio $P_i$ is defined by $(f_i^{\text{firm}}, f_i^{\text{renew}}, f_i^{\text{stor}})$ with $f_i^{\text{firm}} + f_i^{\text{renew}} + f_i^{\text{stor}} = 1$.

The feasible simplex is discretized at 10\% increments subject to $f^{\text{firm}} \in [0.20, 0.70]$, $f^{\text{renew}} \in [0.10, 0.60]$, $f^{\text{stor}} \in [0.10, 0.30]$, yielding $k = 16$ portfolios (Appendix~\ref{app:portfolios}). This systematic approach avoids selection bias and ensures uniform coverage of the design space.

\subsection{Performance Measure}

For portfolio $P_i$ in replication $j$:
\begin{equation}\label{eq:Y}
  Y_{ij} = F_i + C_{ij} + \lambda \, U_{ij},
\end{equation}
where $F_i$ is the annualized fixed cost (deterministic), $C_{ij}$ is simulated annual operating cost, $U_{ij}$ is unserved energy (MWh), and $\lambda = \$10{,}000$/MWh is the value of lost load. Operating cost is calculated from the hourly dispatch simulation:
\begin{equation}\label{eq:C}
  C_{ij} = \sum_{t=1}^{8760} \sum_{g \in \mathcal{G}} c_g \, x_{g,i,j,t},
\end{equation}
where $c_g$ is marginal cost and $x_{g,i,j,t}$ is dispatched output in hour $t$. Both $C_{ij}$ and $U_{ij}$ are a result from the simulation described in Section~\ref{sec:model}; they are not sampled. The objective is $i^* = \arg\min_i \, \theta_i$ where $\theta_i = \mathbb{E}[Y_i]$.

\subsection{Research Questions}

\begin{enumerate}[leftmargin=*, itemsep=2pt]
  \item Within the 12~GW budget and the original 16-portfolio simplex, which mix minimizes expected total cost under uncertainty and does this differ from the deterministic recommendation, and under what structural conditions (capacity budget, simplex bounds) does that difference hold?
  \item Among the leading candidates, which uncertainty factors move the mean cost, and do any tested factor levels change the ranking? (This is a conditional sensitivity analysis on the top-4, and doesn't do the reshuffling test over the full simplex.)
\end{enumerate}

%==========================================================================
% SECTION 3: MODEL DESCRIPTION
%==========================================================================
\section{Model Description}\label{sec:model}

\subsection{System Architecture}

Each portfolio allocates capacity across four firm sub-types, two renewable technologies, and battery storage. Firm generators are dispatched in merit order (Table~\ref{tab:merit}); renewables have zero marginal cost; storage follows a rule-based logic. Full parameter values are in Appendix~\ref{app:params}.

\begin{table}[h]
\centering
\caption{Firm generator merit-order stack.}\label{tab:merit}
\begin{tabular}{lcccc}
\toprule
Sub-type & Share & MC (\$/MWh) & FOR & Capital (\$/kW-yr) \\
\midrule
Nuclear       & 20\% & 22 & 5\% & 400 \\
Flexible firm & 35\% & 35 & 6\% & 100 \\
Mid-merit gas & 30\% & 50 & 5\% & 70 \\
Gas peaker    & 15\% & 80 & 8\% & 70 \\
\bottomrule
\end{tabular}
\end{table}

Nuclear marginal cost of \$22/MWh reflects the variable operating cost of refurbished Ontario nuclear units (fuel plus variable O\&M), not the all-in levelized cost of electricity.

\subsection{Hourly Dispatch Algorithm}

Each replication simulates 8,760 hours of system operation for a given portfolio. The dispatch model uses a sequential priority structure: renewable output is absorbed first (zero marginal cost), storage reacts to net-load conditions, and firm generators are dispatched in ascending marginal-cost order to serve any remaining demand. This merit-order approach reflects standard practice in power system operations, where the cheapest available generation is dispatched first to minimize operating cost. The model does not co-optimize storage and firm dispatch; this simplification keeps the project manageable while preserving the core cost and adequacy dynamics that differentiate portfolios.

For each hour $t$, the dispatch proceeds as follows:

\begin{enumerate}[leftmargin=*, itemsep=2pt]
  \item \textbf{Demand.} $D(t) = D_{\text{base}}(t)(1 + \sigma_d \varepsilon(t))$, $\varepsilon(t) \sim N(0,1)$, $\sigma_d = 0.09$. Multiplicative noise lets absolute deviations scale with load. Demand noise is the only input dimension that does not receive NORTA AR(1) treatment in this study, which is an internal asymmetry with the correlated renewable inputs; the consequences are discussed in Section~\ref{sec:conclusion}.
  \item \textbf{Renewables.} Wind and solar capacity factors are drawn from NORTA AR(1) Beta distributions (Section~\ref{sec:inputs}). Aggregate renewable output is $R(t) = \text{MW}_w \cdot \text{CF}_w(t) + \text{MW}_s \cdot \text{CF}_s(t)$.
  \item \textbf{Residual demand.} $D_{\text{res}}(t) = \max(0,\; D(t) - R(t))$. This is the load that must be served by storage and firm generation.
  \item \textbf{Storage.} The battery discharges when residual demand exceeds the 70th percentile of the base demand profile and charges from surplus when residual demand falls below the 30th percentile, subject to 85\% round-trip efficiency and state-of-charge, power, and remaining-demand limits. These are stylized threshold rules, not co-optimized dispatch; the simplification is noted in Section~\ref{sec:conclusion}.
  \item \textbf{Firm dispatch.} Available firm generators are dispatched in merit order (Table~\ref{tab:merit}). For each sub-type $g$, dispatched output is $x_g(t) = \min(\text{cap}_g \cdot a_g(t),\; D_{\text{rem}}(t))$, where $a_g(t)$ is the outage-adjusted availability fraction. The algorithm walks up the merit-order stack, accumulating operating cost as $x_g(t) \cdot c_g$ and subtracting dispatched output from remaining demand.
  \item \textbf{Unserved energy.} $u(t) = \max(0,\; D_{\text{rem}}(t))$ after all dispatch. Any demand that cannot be served after exhausting renewables, storage, and all available firm generation is recorded as unserved energy for that hour.
\end{enumerate}

\paragraph{Forced outage modeling.}
Each firm sub-type is represented by 10 discrete units of equal size. For each sub-type on each day, each unit independently draws a Bernoulli$(1 - \text{FOR})$ availability indicator, producing discrete availability steps of $0\%, 10\%, \ldots, 100\%$. Outage states persist for 24 hours. This per-unit, per-day granularity is coarser than a continuous-time Markov model but captures the daily availability risk that drives adequacy outcomes.

\vspace{0.5em}
Annual outputs are accumulated over all 8,760 hours: $C_{ij} = \sum_t \sum_g c_g x_{g,t}$, \; $U_{ij} = \sum_t u(t)$, \; and \; $Y_{ij} = F_i + C_{ij} + \lambda U_{ij}$.

\subsection{Input Models and Empirical Calibration}\label{sec:inputs}

Input distributions are calibrated from IESO hourly Ontario demand \citep{IESO2025} and CODERS hourly capacity factors for 2,984 Ontario generation sites \citep{CODERS2025}.

\paragraph{Demand.}
The base profile $D_{\text{base}}(t)$ is fitted via OLS to 2024 IESO data using two seasonal harmonics (annual and semi-annual, capturing Ontario's winter/summer demand peaks) plus a diurnal harmonic, achieving $R^2 = 0.59$. The noise parameter $\sigma_d = 0.09$ is the standard deviation of relative residuals $(D_{\text{obs}} - D_{\text{base}})/D_{\text{base}}$.

\paragraph{Wind.}
Hourly aggregate CFs are fitted to Beta$(6.73, 15.69)$ via MLE (KS $= 0.021$, mean $= 0.300$). To preserve the strong temporal persistence (empirical lag-1 ACF $\hat\rho = 0.995$), CFs are generated via NORTA with AR(1) dependence:
\begin{equation}\label{eq:norta}
  Z(t) = \rho Z(t{-}1) + \sqrt{1-\rho^2}\,\epsilon(t), \quad
  \text{CF}(t) = F_{\text{Beta}}^{-1}(\Phi(Z(t))).
\end{equation}
This preserves the Beta marginal while inducing hourly autocorrelation. The distinction from iid sampling is consequential: iid draws overestimate hourly variability while underestimating sustained low-output periods that drive adequacy risk (Appendix~\ref{app:inputs}).

\paragraph{Solar.}
$\text{CF}_s(t) = E(t) \cdot K(t)$, where $E(t)$ is a sinusoidal diurnal envelope (hours 7--21) and $K(t) \sim \text{Beta}(1.77, 6.91)$ via NORTA AR(1) with $\rho = 0.928$, calibrated from CODERS.

\paragraph{Forced outages.}
See Section~\ref{sec:model}: 10 units per sub-type, independent daily Bernoulli$(1 - \text{FOR})$ draws.

%==========================================================================
% SECTION 4: EXPERIMENTAL DESIGN
%==========================================================================
\section{Experimental Design}\label{sec:design}

Each simulation run is a terminating simulation over 8,760 hours (one operating year), so output analysis uses independent replications rather than batch means. Error is controlled at three levels: the Kim--Nelson ranking-and-selection procedure guarantees $P(\text{CS}) \geq 0.95$ for the indifference-zone parameter $\delta = \$50$M, the top-4 evaluation uses Bonferroni-corrected pairwise 95\% confidence intervals with $\alpha_{\text{pair}} = 0.05/6 = 0.0083$, and relative half-widths are reported for every mean (Table~\ref{tab:topk}).

\subsection{Common Random Numbers}

All portfolios share identical random inputs within each replication, implemented via NumPy's \texttt{SeedSequence} spawning. Three independent streams (demand, outages, renewables) are generated deterministically from a base seed and replication index. Under CRN, positive covariance between outputs reduces $\text{Var}(Y_i - Y_\ell)$, which is the quantity that drives pairwise inference in ranking and selection. The empirical variance reduction on pairwise differences ranges from 2$\times$ to 35$\times$ (Appendix, Figure~\ref{fig:crn}).

\subsection{Kim--Nelson Ranking and Selection}

The Kim--Nelson (KN) sequential ranking-and-selection procedure \citep{KimNelson2001} selects the portfolio with smallest $\theta_i$ with probability of correct selection at least $1 - \alpha$, provided the best system is at least $\delta$ better than the second best. Parameters used throughout are $\alpha = 0.05$, $\delta = \$50$M (about 1\% of annual system cost), and $n_0 = 30$. KN begins with $n_0$ CRN replications on all $k = 16$ portfolios, computes pairwise difference variances $S^2_{i\ell}$ (which absorb the CRN covariance automatically), and at each stage $r \geq n_0$ eliminates portfolio $i$ whenever some surviving competitor $\ell$ satisfies $\bar{Y}_i(r) - \bar{Y}_\ell(r) > W_{i\ell}(r)$, with continuation region
\begin{equation}\label{eq:kn-threshold}
  W_{i\ell}(r) = \max\!\left\{0,\; \frac{\delta}{2r}\!\left(\frac{h^2\, S^2_{i\ell}}{\delta^2} - r\right)\right\}
\end{equation}
with $h^2 = 2\eta(n_0 - 1)$ and $\eta$ the Bonferroni-corrected continuation constant for $k$ systems at level $\alpha$. The region $W_{i\ell}(r)$ shrinks with $r$ and reaches zero in finite time, guaranteeing termination with one survivor.

\subsection{Deterministic Benchmark Gap}

The deterministic recommendation $P_{\text{DET}}$ is obtained by setting all inputs to their expected values and selecting the portfolio with the lowest $Y$ within the original simplex. $P_{\text{DET}}$ and the KN-selected portfolio $P_{\text{KN}}$ are then evaluated over 200 paired CRN replications. The DBG is defined as $\bar{Y}_{\text{DET}} - \bar{Y}_{\text{KN}}$ and, as used in this paper, is a scenario-specific comparison within the 12~GW budget and the original simplex and should not be defined as a universal planning penalty or a general statement about the value of uncertainty-aware planning. Section~\ref{sec:extensions} re-runs the same construction at higher capacity levels and reports how the DBG behaves outside the 12~GW scenario.

\subsection{Sensitivity Analysis}

A $2^3$ factorial tests demand noise ($\sigma_d \in \{0.03, 0.15\}$), forced outage rates (all 1\% vs.\ realistic), and renewable variability (fixed mean vs.\ stochastic). The design is restricted to the top-4 portfolios, with 30 CRN replications per cell.

\subsection{Extensions}

Three additional analyses extend the core evaluation. First, a \textit{stochastic adequacy search} finds, for each top-4 mix, the minimum total capacity $C^*$ such that $\mathbb{E}[\text{EUE}(C)] \leq 10$~MWh/year. This is a one-dimensional simulation-optimization problem (stochastic root finding for $\mathbb{E}[\text{EUE}(C)] - 10 = 0$): at each candidate $C$, 50 CRN replications estimate the noisy objective, and bisection iterates until the threshold is crossed within tolerance. Second, a \textit{UCAP-normalized comparison} re-evaluates six portfolios (spanning 20--70\% firm) each sized to equal effective capacity (8,000~MW UCAP) using representative ELCC values (Appendix~\ref{app:elcc}); this tests whether the equal-installed-MW design biases results toward firm-heavy portfolios. Third, a \textit{post-processing VOLL sensitivity} recomputes $Y = F + C + \lambda U$ over $\lambda \in \{5{,}000, 10{,}000, 15{,}000\}$ using the replication-level components from the extended CRN run.

%==========================================================================
% SECTION 5: RESULTS
%==========================================================================
\section{Results}\label{sec:results}

\subsection{In-Budget Portfolio Selection Within the Original Simplex}

The main experiment ranks the 16 portfolios of the 10\%-step simplex under a fixed 12~GW budget. Both the budget and the 70\% firm-share ceiling are binding design choices: the 12~GW figure is the stylized stress scenario introduced in Section~\ref{sec:intro}, and the firm-share ceiling curtails the design space directly. Section~\ref{sec:extensions} tests the sensitivity of the main result to each of these two constraints. The result below is the in-budget winner within that constrained design. KN selects P04 (70/20/10) after 77 replications (Figure~\ref{fig:kn}). Fourteen of the sixteen portfolios are eliminated at the initial stage ($n_0 = 30$), leaving P04 and P01 (40/50/10); P01 is then eliminated at stage~77. Table~\ref{tab:topk} reports the extended evaluation of the top-4 survivors over 200 CRN replications.

\begin{figure}[ht]
\centering
\includegraphics[width=0.78\textwidth]{figures/fig2_kn_elimination.png}
\caption{KN elimination timeline. Fourteen of sixteen portfolios are eliminated at $n_0 = 30$; P01 is eliminated at stage 77, leaving P04 (70/20/10) as the in-budget winner within the original simplex.}\label{fig:kn}
\end{figure}

\begin{table}[ht]
\centering
\caption{Top-4 portfolio evaluation (200 replications). $\bar{U}$ is mean annual unserved energy in MWh; $\lambda = \$10{,}000$/MWh. RHW is the relative half-width of the 95\% confidence interval for mean $Y$.}\label{tab:topk}
\small
\begin{tabular}{lcccccc}
\toprule
Portfolio & Mean $Y$ (\$M) & $F_i$ (\$M) & $\bar{C}$ (\$M) & $\bar{U}$ (MWh) & $\lambda\bar{U}$ (\$M) & RHW \\
\midrule
P04 (70/20/10) & 4,306 & 1,694 & 1,835 &  77{,}700 & 777   & 0.38\% \\
P10 (70/10/20) & 5,060 & 1,732 & 1,874 & 145{,}400 & 1,454 & 0.40\% \\
P03 (60/30/10) & 8,272 & 1,660 & 1,846 & 476{,}600 & 4,766 & 0.83\% \\
P09 (60/20/20) & 10,785 & 1,698 & 1,882 & 720{,}500 & 7,205 & 0.64\% \\
\bottomrule
\end{tabular}
\end{table}
The decomposition in Table~\ref{tab:topk} shows where the gap comes from. Operating costs vary by less than 3\% across the top-4, but the adequacy penalty ranges from \$777M for P04 to \$7{,}205M for P09, so the gap between these mixes is driven almost entirely by the adequacy penalty. Low firm capacity interacts with the temporally correlated wind droughts produced by the NORTA AR(1) model, and those sustained low-output episodes, which do not appear under iid sampling or mean-input evaluation, generate the heavy right tails visible in Figure~\ref{fig:box} for P03 and P09.

All six pairwise differences among the top-4 are statistically significant under Bonferroni correction ($\alpha_{\text{pair}} = 0.0083$), and all relative half-widths fall below 1\%. KN terminates after 77 replications, well under the 500-replication cap, which is consistent with the variance reduction CRN achieves on pairwise differences (Figure~\ref{fig:crn}).


\subsection{Deterministic Benchmark Gap (12~GW Budget)}\label{sec:dbg}

Fixing all inputs at their expected values (zero demand noise, no forced outages, annual-mean renewables) makes P03 (60/30/10) the cheapest portfolio within the original simplex at \$3{,}503M. That choice differs from the in-budget stochastic selection P04, and the gap is
\begin{equation}\label{eq:dbg}
  \text{DBG} = \bar{Y}_{\text{P03}} - \bar{Y}_{\text{P04}} = \$3{,}968\text{M}, \quad 95\%\;\text{CI:}\; [\$3{,}905,\; \$4{,}031]\,\text{M}.
\end{equation}
Under the 12~GW budget, mean-input evaluation selects a portfolio whose adequacy deteriorates sharply once the tail events are simulated (Figure~\ref{fig:dbg}). The arithmetic behind this is straightforward. P03 saves \$34M in annualized fixed cost relative to P04 because its lower firm fraction needs less capital-intensive generation, but that saving only holds if residual demand never exceeds the available firm capacity. With stochastic inputs, demand noise ($\sigma_d = 0.09$), forced outages (5--8\% FOR), and NORTA-modelled wind droughts push residual demand above P03's 7{,}200~MW of outage-reduced firm capacity during sustained low-wind episodes, and the adequacy penalty more than erases the fixed-cost saving.

The \$3{,}968M figure is specific to this setup: the 12~GW budget, the original simplex, and the stochastic input structure used here. It is not a universal planning penalty and it should not be read as a general statement about uncertainty-aware planning. Section~\ref{sec:extensions} runs the same construction at higher capacities and finds that the gap closes within the indifference zone, which answers Research Question~1 conditional on the 12~GW scenario and the original simplex bounds.

\begin{figure}[ht]
\centering
\includegraphics[width=0.60\textwidth]{figures/fig6_dbg.png}
\caption{Deterministic Benchmark Gap. Under mean-input assumptions, P03 appears cheaper. Under stochastic evaluation with 200 paired CRN replications, P04 costs \$3,968M less per year.}\label{fig:dbg}
\end{figure}

\subsection{Conditional Factorial Sensitivity (Top-4)}

A $2^3$ factorial on the top-4 KN survivors is run with 30 CRN replications per cell. This is a conditional sensitivity analysis among leading candidates: the twelve eliminated portfolios are not re-evaluated under the factorial factor levels, so the factorial cannot speak to whether uncertainty assumptions would reshuffle the full 16-portfolio ranking. Within the top-4 set, demand noise (\$2,883M main effect) and forced outages (\$2,838M) dominate, while renewable variability contributes only \$117M (Figure~\ref{fig:tornado}). The small renewable effect reflects the concentrated Beta(6.73, 15.69) marginal calibrated from fleet-level CODERS data.

The negative demand$\times$outage interaction ($-$\$699M) reflects the fact that once either factor alone is at its high level, inadequate hours are already experiencing unserved energy. The second stressor largely hits the same hours, so it adds less unserved energy than it would on its own. Across all eight configurations, P04 remains the in-budget winner among the top-4. Research Question~2 therefore has a conditional answer. Demand noise and forced outages move the \emph{magnitude} of costs among these four portfolios but do not reshuffle their ranking within the 12~GW budget and the original simplex. Whether the identity of the winner survives changes in the capacity budget or the simplex boundary is a stronger question, and Section~\ref{sec:extensions} shows that it does not.

\begin{figure}[ht]
\centering
\includegraphics[width=0.80\textwidth]{figures/fig4_tornado.png}
\caption{Conditional factorial sensitivity on the top-4 portfolios only. Demand noise and forced outages dominate; renewable variability has minimal impact. P04 remains the in-budget winner across all eight factorial configurations within this restricted set.}\label{fig:tornado}
\end{figure}

\subsection{VOLL Sensitivity}

Re-ranking the top-4 for $\lambda \in \{5{,}000,\,10{,}000,\,15{,}000\}$~\$/MWh by post-processing the existing $(F, C, U)$ outputs leaves P04 as the in-budget winner of the original simplex at all three values, so the 12~GW selection is not an artefact of a single VOLL calibration point within this range.

\subsection{Adequacy Search}\label{subsec:adequacy}

Treating total capacity as a one-dimensional decision variable and bisecting on $\mathbb{E}[\text{EUE}(C)] \leq 10$~MWh/year per portfolio, P04 reaches the adequacy threshold at 15.8~GW, P03 at 17.5~GW, and P09 at 18.2~GW. Among these three, each evaluated at its own adequacy point, the cheapest is P04 at \$3{,}763M. This is a per-portfolio construction: every mix is compared at a different total capacity. Section~\ref{sec:extensions} runs the complementary experiment in which the same top-4 portfolios are evaluated at a common fixed capacity (12, 15.8, and 18~GW), which is what allows a direct between-portfolio ranking at each level. The two results answer different questions and are not in tension.

\subsection{UCAP-Normalized Comparison and Empirical ELCC}

One reasonable objection to the main result is that the equal-installed-MW design favours firm-heavy mixes. To check this, six portfolios spanning 20\% to 70\% firm are re-evaluated at equal UCAP (8{,}000~MW). P04 is still cheapest (\$9{,}186M against \$10{,}167M for P03), which suggests that static ELCC, a peak-period metric, does not capture the sustained low-renewable events that drive the adequacy penalty in this simulator (Appendix~\ref{app:results}, Figure~\ref{fig:ucap}).

A separate empirical ELCC experiment (Appendix~\ref{app:elcc_emp}) attempts to quantify the gap. For each technology, 1{,}000~MW is added to a stressed 5{,}000~MW firm-only baseline and the resulting reduction in unserved energy is normalized against a perfect-firm reference. The implied values are wind 0.25 (vs.\ 0.18 static), solar 0.07 (vs.\ 0.35), and storage 0.13 (vs.\ 0.50). These are properties of the stressed baseline and should not be read as the true ELCC of Ontario wind, solar, or storage at the 12~GW portfolio size. For wind and solar, the direction (static values overstate the adequacy contribution of intermittent resources once sustained low-renewable events bind) is consistent with the UCAP-normalized comparison. The storage result is weaker. Storage dispatch here is a threshold heuristic rather than a co-optimized program, which entangles the empirical storage ELCC with the heuristic itself. A more sophisticated dispatch model would very likely narrow the gap with the static value, and no claim about the true adequacy value of storage should be drawn from the present implementation.

\subsection{Structural Sensitivity: Boundary and Capacity Probes}\label{sec:extensions}

Two structural features of the main experiment warrant a direct test. P04 sits at the upper firm-share boundary of the simplex, and the 12~GW budget itself is inadequate by Table~\ref{tab:topk} and by the adequacy search. Both are design choices rather than modelling details, and the main result should be reported in a way that makes clear how it depends on them. Appendix~\ref{app:structural} contains the details of the two probes described below.

\paragraph{Boundary probe.} Four portfolios with firm share in $\{0.75, 0.80, 0.85\}$ are constructed outside the original simplex and re-evaluated at 12~GW with 100 CRN replications each. Mean $Y$ falls monotonically as firm share increases, and P$_{b3}$ (85/10/05) reaches \$3{,}553M, which is \$746M cheaper than P04's \$4{,}299M. P04 is not the true in-budget optimum under the 12~GW budget; it was simply the cheapest mix the original simplex could reach.

\paragraph{Capacity sweep.} The top-4 portfolios (P04, P10, P03, P09) are re-evaluated at 12~GW, 15.8~GW, and 18~GW. At 12~GW, P04 is the cheapest of the four. At 15.8~GW and 18~GW, P03 (60/30/10) is the cheapest of the four, beating P04 by \$68M and \$113M respectively. At both higher capacities, the deterministic and stochastic recommendations agree on P03 within this set, and the DBG for the set closes within the indifference zone. This is a within-set result on the four portfolios in the sweep, not a claim over an expanded design space, but it is enough to show that the \$3{,}968M DBG from \S\ref{sec:dbg} is tied to the inadequate 12~GW scenario rather than to uncertainty-aware evaluation in general.

The two tests sharpen the main result. In the budget-constrained scenario the system is tail-risk dominated, firm capacity wins, and the in-budget optimum sits beyond the 70\% boundary. Once the budget is relaxed to adequacy, the system becomes cost-dominated, more renewable-heavy mixes become cheaper within the set re-evaluated, and the two evaluation methods converge on that set. Two further internal checks (an OCBA consistency note and a $\pm 30\%$ capital-cost sensitivity) and an internal verification suite appear in Appendices~\ref{app:vv}, \ref{app:capsens}, and \ref{app:ocba}.

%==========================================================================
% SECTION 6: CONCLUSION
%==========================================================================
\section{Conclusion}\label{sec:conclusion}

In a stylized one-zone system with a fixed 12~GW capacity budget, simulation-based uncertainty-aware evaluation meaningfully changes the planner's recommendation only in the setting where the budget binds. Within the 12~GW budget and the original 16-portfolio simplex, stochastic evaluation picks P04 (70/20/10) and mean-input evaluation picks P03 (60/30/10), giving a Deterministic Benchmark Gap of \$3{,}968M per year. When the same four portfolios are re-evaluated at 15.8~GW and 18~GW, both methods agree on P03 within that set and the DBG closes within the indifference zone. The \$3{,}968M figure is specific to the 12~GW stress scenario and the original simplex; it is not a general planning penalty. A separate boundary test shows that relaxing the 70\% firm-share ceiling exposes a cheaper 12~GW mix, P$_{b3}$ (85/10/05), which beats P04 by \$746M. P04 is the in-budget winner within the original constrained design, not an absolute optimum.

Several design choices in the methodology matter substantively. The NORTA AR(1) wind model preserves both the empirical Beta marginal and the temporal autocorrelation ($\rho = 0.995$) of Ontario wind data. This matters because iid sampling would understate the sustained low-output episodes that drive the P04--P03 gap. CRN reduces pairwise-difference variance by a factor of 2 to 35, which is what allows KN to terminate early. The adequacy search and the capacity sweep are complementary: the first solves a one-dimensional stochastic root-finding problem on per-portfolio adequacy, and the second re-runs the same four portfolios at common fixed capacities. Together they expose the scenario dependence of the main result. These findings are consistent with the broad motivation of \citet{ParkBaldick2016}, but the methodology and decision problem differ enough that the present study should not be read as a methodological adaptation of that paper.

Several simplifications remain. The most consequential is the treatment of demand noise as iid Gaussian around an OLS harmonic fit. Demand is simultaneously the largest single factor in the top-4 factorial (main effect \$2{,}883M) and the only input dimension that does not receive NORTA AR(1) treatment. This is an asymmetry that no direct experiment here has bounded. Whether autocorrelated demand residuals would widen or reverse the 12~GW DBG is not a question this study can answer, and it is flagged as a limitation rather than as a prediction. Storage dispatch is a threshold heuristic, which entangles the empirical storage ELCC result with the dispatch rule itself; no claim about the true adequacy value of storage should be taken from the present implementation. NORTA AR(1) approximates normal-space $\rho$ by the target autocorrelation, with negligible error at $\rho = 0.995$, and forced outages are represented at a 10-unit per-sub-type discretization that is coarser than a continuous-time Markov model. These components are verified only against their own analytical targets (Appendix~\ref{app:vv}), so the study claims internal consistency but not external validation against a ground-truth system. The natural extension is a full stochastic capacity-expansion formulation that jointly treats the capacity budget and the resource mix. Given the scenario dependence observed here, the value of uncertainty-aware evaluation in that fuller problem is likely concentrated where the budget constraint binds.

%==========================================================================
% REFERENCES
%==========================================================================
\begingroup
\scriptsize
\setlength{\bibsep}{1.5pt plus 0.2ex}
\bibliographystyle{abbrvnat}
\bibliography{references}
\endgroup


%==========================================================================
% APPENDIX
%==========================================================================
\appendix

\section{Consolidated Parameters}\label{app:params}

\begin{tabular}{lll}
\toprule
Parameter & Value & Basis \\
\midrule
\multicolumn{3}{l}{\textit{System}} \\
Peak demand & 8,000 MW & Stylized \\
Reserve margin & 50\% & Accounts for renewable ELCC \\
Total capacity & 12,000 MW & Peak $\times$ 1.5 \\
\midrule
\multicolumn{3}{l}{\textit{Demand}} \\
Profile & 2-harmonic OLS & IESO 2024 \citep{IESO2025} ($R^2 = 0.59$) \\
Noise $\sigma_d$ & 0.09 & Residual std from fit \\
\midrule
\multicolumn{3}{l}{\textit{Renewables}} \\
Wind CF & Beta(6.73, 15.69) & MLE, CODERS \citep{CODERS2025} \\
Wind $\rho$ & 0.995 & Lag-1 ACF, CODERS \\
Solar cloud & Beta(1.77, 6.91) & MLE, CODERS \\
Solar $\rho$ & 0.928 & Lag-1 ACF, CODERS \\
Daylight hours & 07:00--21:00 & CODERS (UTC$-$5) \\
Wind/solar split & 60/40 & Stylized \\
\midrule
\multicolumn{3}{l}{\textit{Storage}} \\
Duration & 4 hours & Industry standard \\
Round-trip eff.\ & 85\% & NREL ATB \citep{NRELATB2023} \\
Discharge/charge & 70th/30th pctile & Stylized rules \\
\midrule
\multicolumn{3}{l}{\textit{Cost}$^*$} \\
VOLL ($\lambda$) & \$10,000/MWh & Mid-range \\
Nuclear capital & \$400/kW-yr & High overnight, 30-yr \\
Flex firm capital & \$100/kW-yr & Moderate \\
Mid-merit capital & \$70/kW-yr & Combined cycle \\
Peaker capital & \$70/kW-yr & Simple cycle \\
Wind capital & \$130/kW-yr & Onshore \\
Solar capital & \$100/kW-yr & Utility-scale \\
Storage capital & \$150/kW-yr & Battery, 15-yr life \\
\midrule
\multicolumn{3}{l}{\textit{KN procedure}} \\
$\alpha$ / $\delta$ / $n_0$ & 0.05 / \$50M / 30 & \\
\bottomrule
\multicolumn{3}{l}{\footnotesize $^*$Capital costs are stylized, informed by NREL ATB \citep{NRELATB2023}.} \\
\end{tabular}


\section{Portfolio Definitions}\label{app:portfolios}

\begin{table}[H]
\centering
\caption{All 16 candidate portfolios (12,000 MW total).}\label{tab:allport}
\small
\begin{tabular}{clrrrr}
\toprule
ID & Label & Firm & Renew & Storage & $F_i$ (\$M) \\
\midrule
P00 & 30/60/10 & 3,600 & 7,200 & 1,200 & 1,557 \\
P01 & 40/50/10 & 4,800 & 6,000 & 1,200 & 1,591 \\
P02 & 50/40/10 & 6,000 & 4,800 & 1,200 & 1,625 \\
P03 & 60/30/10 & 7,200 & 3,600 & 1,200 & 1,660 \\
P04 & 70/20/10 & 8,400 & 2,400 & 1,200 & 1,694 \\
P05 & 20/60/20 & 2,400 & 7,200 & 2,400 & 1,561 \\
P06 & 30/50/20 & 3,600 & 6,000 & 2,400 & 1,595 \\
P07 & 40/40/20 & 4,800 & 4,800 & 2,400 & 1,630 \\
P08 & 50/30/20 & 6,000 & 3,600 & 2,400 & 1,664 \\
P09 & 60/20/20 & 7,200 & 2,400 & 2,400 & 1,698 \\
P10 & 70/10/20 & 8,400 & 1,200 & 2,400 & 1,732 \\
P11 & 20/50/30 & 2,400 & 6,000 & 3,600 & 1,600 \\
P12 & 30/40/30 & 3,600 & 4,800 & 3,600 & 1,634 \\
P13 & 40/30/30 & 4,800 & 3,600 & 3,600 & 1,668 \\
P14 & 50/20/30 & 6,000 & 2,400 & 3,600 & 1,702 \\
P15 & 60/10/30 & 7,200 & 1,200 & 3,600 & 1,736 \\
\bottomrule
\end{tabular}
\end{table}


\section{ELCC Values}\label{app:elcc}

\begin{table}[H]
\centering
\caption{ELCC values for UCAP normalization.}\label{tab:elcc}
\small
\begin{tabular}{lcc}
\toprule
Technology & ELCC & Basis \\
\midrule
Nuclear      & 94\% & High availability baseload \\
Flexible firm & 92\% & Gas/hydro, slight derates \\
Mid-merit gas & 95\% & CCGT \\
Gas peaker   & 92\% & Cycling stress \\
Wind         & 18\% & IESO capacity auction \\
Solar        & 35\% & Summer peak contribution \\
Storage (4h) & 50\% & Duration-limited \\
\bottomrule
\end{tabular}
\end{table}


\section{Input Fitting Diagnostics}\label{app:inputs}

\begin{figure}[H]
\centering
\includegraphics[width=0.80\textwidth]{figures/input_demand_fit.png}
\caption{Demand profile OLS fit to IESO 2024 hourly Ontario demand. Left: observed vs.\ fitted. Right: residual distribution ($\sigma_d = 0.09$).}\label{fig:demand_fit}
\end{figure}

\begin{figure}[H]
\centering
\includegraphics[width=0.80\textwidth]{figures/input_wind_fit.png}
\caption{Wind CF Beta(6.73, 15.69) MLE fit to CODERS Ontario. Left: empirical vs.\ fitted PDF. Right: Q--Q plot (KS $= 0.021$).}\label{fig:wind_fit}
\end{figure}

\begin{figure}[H]
\centering
\includegraphics[width=0.80\textwidth]{figures/input_solar_fit.png}
\caption{Solar CF: diurnal envelope $\times$ Beta(1.77, 6.91), calibrated from CODERS Ontario.}\label{fig:solar_fit}
\end{figure}


\section{Supplementary Results}\label{app:results}

\begin{figure}[H]
\centering
\includegraphics[width=0.70\textwidth]{figures/fig3_y_distributions.png}
\caption{Distribution of $Y_{ij}$ for top-4 portfolios. P03 and P09 exhibit heavy right tails from high-penalty replications during sustained low-renewable periods.}\label{fig:box}
\end{figure}

\begin{figure}[H]
\centering
\includegraphics[width=0.85\textwidth]{figures/fig7_crn_efficiency.png}
\caption{CRN variance reduction: pairwise difference variance under independent vs.\ CRN sampling. Efficiency ranges from 2$\times$ (P04 vs.\ P09) to 35$\times$ (P03 vs.\ P09).}\label{fig:crn}
\end{figure}

\begin{figure}[H]
\centering
\includegraphics[width=0.90\textwidth]{figures/fig12_ucap_comparison.png}
\caption{UCAP-normalized comparison at equal effective capacity (8~GW UCAP). Left: mean total cost per mix; right: cost decomposition. P04 remains cheapest after normalization.}\label{fig:ucap}
\end{figure}


\section{Internal Verification Suite}\label{app:vv}

\texttt{src/verification.py}. Five automated internal-consistency checks, all pass. These are internal verification, not external validation: they show that the simulator behaves as specified and that its stochastic machinery reproduces its own analytical targets. They do not validate the stylized one-zone model against a ground-truth system.

\begin{table}[H]
\centering
\small
\begin{tabular}{p{4.2cm}p{5.0cm}p{4.6cm}c}
\toprule
Check & Observed (simulator) & Expected (analytical) & Pass \\
\midrule
Deterministic limit of P04 (zero noise, zero FOR)
    & $U_{ij}=0$;\; $Y_{ij}=\$3{,}505$M
    & $U = 0$;\; $Y = F + C_{\text{det}}$
    & \checkmark \\[2pt]
Input mean recovery (50 reps)
    & demand 6{,}549~MW (${<}0.01\%$);\newline wind CF 0.304 (1.2\%)
    & demand 6{,}549~MW;\newline Beta marginal 0.300
    & \checkmark \\[2pt]
Outage-rate recovery (50 reps, all four sub-types)
    & within 0.1 pp of $1 - \text{FOR}$
    & $1 - \text{FOR}$
    & \checkmark \\[2pt]
Monotonicity: low-firm vs.\ high-firm EUE
    & P05: 21.2M MWh\newline P04: 0.07M MWh (285$\times$)
    & $\text{EUE}_{\text{low firm}} > \text{EUE}_{\text{high firm}}$
    & \checkmark \\[2pt]
CRN variance reduction on P04 $-$ P10
    & $\operatorname{sd} = \$63.7$M (8.8$\times$ reduction)
    & $<$ independent $\operatorname{sd} = \$189.1$M
    & \checkmark \\
\bottomrule
\end{tabular}
\end{table}

\noindent The rule-based storage and discrete 10-unit outage model (Section~\ref{sec:conclusion}) are reported limitations rather than V\&V failures, because neither has an analytical target to verify against.


\section{Capital Cost Sensitivity}\label{app:capsens}

\texttt{src/capital\_sensitivity.py}. One-at-a-time $\pm 30\%$ sweep on each of the seven capital-cost parameters (Appendix~\ref{app:params}); $F_i$ recomputed, $C_{ij}$ and $U_{ij}$ held fixed from the 200-replication CRN baseline. P04 remains best in all 35 perturbed configurations; minimum P10$-$P04 Mean $Y$ gap $= \$700$M (storage $+30\%$), well above the \$50M indifference zone.

\begin{figure}[H]
\centering
\includegraphics[width=0.72\textwidth]{figures/fig_capital_sensitivity.png}
\caption{P10$-$P04 Mean $Y$ gap under $\pm 30\%$ one-at-a-time perturbations of each capital cost parameter. Dashed red: baseline gap. Dotted grey: \$50M indifference zone.}\label{fig:capsens}
\end{figure}


\section{Empirical ELCC Experiment}\label{app:elcc_emp}

\texttt{src/empirical\_elcc.py}. Method: add a 1{,}000~MW dose of technology $X$ to a stressed 5{,}000~MW firm-only baseline $S_0$ (realistic FOR, no renewables, no storage), measure $\mathbb{E}[U_X]$ over 100 CRN replications, and define $\text{ELCC}_X^{\text{emp}} = (\mathbb{E}[U_0] - \mathbb{E}[U_X])\,/\,(\mathbb{E}[U_0] - \mathbb{E}[U_{\text{perf}}])$, where $U_{\text{perf}}$ is obtained by adding 1{,}000~MW of firm with all FORs overridden to zero. Baseline $\mathbb{E}[U_0] = 16.14$M~MWh; perfect-firm reference reduction $\Delta U_{\text{perf}} = 9.98$M~MWh. The stressed baseline is necessary because EUE at the 12,000~MW capacity of the main experiment is too small to resolve the marginal contribution of a 1{,}000~MW addition at a reasonable replication count. ELCC is load- and portfolio-dependent, so the numerical factors reported here are specific to this stressed configuration and should not be read as the true ELCC of Ontario wind, solar, or storage in the 12,000~MW portfolios. The qualitative conclusion (static ELCC overstates solar and storage contribution when sustained low-renewable events are the binding adequacy constraint) is what transfers.

\begin{figure}[H]
\centering
\includegraphics[width=0.72\textwidth]{figures/fig_elcc_empirical.png}
\caption{Empirical vs.\ static ELCC, 1,000 MW dose on the 5,000 MW stressed firm baseline. Solar and storage come in at ${\sim}1/5$ and ${\sim}1/4$ of their static planning values.}\label{fig:elcc_emp}
\end{figure}


\section{OCBA Consistency Note}\label{app:ocba}

\texttt{src/ocba.py}. OCBA \citep{ChenOCBA2000} was run against KN across five macro-replications at different base seeds purely as a consistency check on the 12~GW / original-simplex KN selection. At a matched simulation-call budget (mean 528 calls), OCBA selects P04 in 5/5 runs, matching KN. Five macro-reps are too few for a meaningful $P(\text{CS})$ estimate and the large gap between P04 and its runner-up within the original simplex makes selection easy for either method, so this is \emph{not} a powered procedural comparison and no performance claim between KN and OCBA is made. The check is reported only because an independent R\&S procedure reaching the same conclusion adds confidence that the KN result is not an implementation artefact.


\section{Structural Sensitivity: Boundary and Capacity}\label{app:structural}

\texttt{src/deeper\_sensitivity.py}. Two probes test whether the main 12~GW / 70\%-firm-upper-bound result is a structural artefact.

\begin{figure}[H]
\centering
\begin{minipage}{0.50\textwidth}
  \centering
  \includegraphics[width=\textwidth]{figures/fig_boundary.png}
\end{minipage}\hfill
\begin{minipage}{0.50\textwidth}
  \centering
  \includegraphics[width=\textwidth]{figures/fig_capacity_sweep.png}
\end{minipage}
\caption{Left: boundary probe at 12~GW. P$_{b3}$ (85/10/05) beats P04 by \$746M. Right: capacity sweep; winners ($\star$) flip from P04 at 12~GW to P03 at 15.8~GW and 18~GW.}\label{fig:structural}
\end{figure}

\begin{table}[H]
\centering
\caption{DBG at adequate capacity. Deterministic uses zero demand noise, zero FOR, fixed-mean renewables; stochastic uses 100 CRN replications.}\label{tab:dbg_capacity}
\small
\begin{tabular}{lccc}
\toprule
Total capacity & Deterministic winner & Stochastic winner & DBG (\$M) \\
\midrule
12.0 GW (main)     & P03 (60/30/10) & P04 (70/20/10) & 3{,}968 \\
15.8 GW (P04 adequate) & P03 (60/30/10) & P03 (60/30/10) & 0 \\
18.0 GW (all adequate) & P03 (60/30/10) & P03 (60/30/10) & 0 \\
\bottomrule
\end{tabular}
\end{table}


\section{Code Architecture}\label{app:code}

Python 3.11 (NumPy, SciPy). The codebase comprises sixteen source files totalling approximately 3,000 lines, organized as follows:

\noindent\textbf{Core pipeline:} \texttt{config.py}, \texttt{crn.py}, \texttt{dispatch.py}, \texttt{kn.py}, \texttt{experiment.py}, \texttt{plots.py}, \texttt{main.py}.

\noindent\textbf{Extensions:} \texttt{verification.py}, \texttt{capital\_sensitivity.py}, \texttt{empirical\_elcc.py}, \texttt{ocba.py},\\ \texttt{deeper\_sensitivity.py}, \texttt{adequacy.py}.

The full pipeline runs in approximately 5 minutes on a consumer laptop. Base seed 42 for the core pipeline; seeds 100--104 for OCBA macro-replications. Reproduce via \texttt{python src/main.py} and \texttt{python src/<module>.py} for each extension. All source files are submitted alongside this report on Quercus. 

\vspace{0.5em}
\noindent The full repository is also available at:
\par\noindent\url{https://github.com/Mourssio/capacity-portfolio-eval}.

\section{Generative AI Disclosure}\label{app:ai}

Generative AI (Claude, Anthropic) was used for code development, debugging, statistical methodology review, and drafting support. All modelling decisions, parameter choices, experimental design, interpretation, and final writing are the author's own.

\end{document}
